% Synthetic regression; literal Ukrainian source, never patched by the runner.
\documentclass{article}
\usepackage{fontspec}
\setmainfont{DejaVuSerif.ttf}
\usepackage[english,ukrainian]{babel}
\makeatletter
\typeout{ENGLISH-ID:\number\l@english}
\typeout{UKRAINIAN-ID:\number\l@ukrainian}
\makeatother
\begin{document}
\selectlanguage{ukrainian}
\typeout{UKRAINIAN-ACTIVE:\the\language}
\typeout{UKRAINIAN-CONTENTS:\contentsname}
\typeout{UKRAINIAN-HYPHENS-BEGIN}
% XeTeX kernel showhyphens forces TU/lmr; explicitly select the test font
% inside its diagnostic box, so missing Cyrillic glyphs cannot hide a failure.
\showhyphens{\fontspec{DejaVuSerif.ttf}український українська університет комп'ютерний програмування}
\typeout{UKRAINIAN-HYPHENS-END}
\contentsname: український університет. Іі Її Єє Ґґ.
\selectlanguage{english}
\typeout{ENGLISH-ACTIVE:\the\language}
\typeout{ENGLISH-CONTENTS:\contentsname}
\typeout{ENGLISH-HYPHENS-BEGIN}
\showhyphens{\fontspec{DejaVuSerif.ttf}hyphenation university compatibility}
\typeout{ENGLISH-HYPHENS-END}
\typeout{UKRAINIAN-IN-ENGLISH-CONTROL-BEGIN}
\showhyphens{\fontspec{DejaVuSerif.ttf}український університет}
\typeout{UKRAINIAN-IN-ENGLISH-CONTROL-END}
ASCII I i 123. English retained.
\end{document}
